\section{二级结论卡片}

二级结论是从定义和性质推出来的、使用频率高的结果。
下面每张卡片包含五栏：结论、推导、适用条件、易错点、配套练习。

\begin{table}[htbp]
\centering
\small
\setlength{\tabcolsep}{3pt}
\renewcommand{\arraystretch}{1.5}
\begin{tabular}{p{2.6cm}p{4.5cm}p{1.8cm}p{2.5cm}p{0.9cm}}
\toprule
\key{结论} & \key{推导} & \key{适用条件} & \key{易错点} & \key{层级} \\
\midrule
顶点坐标
$\left(-\dfrac{b}{2a},\dfrac{4ac-b^2}{4a}\right)$
& 对 $ax^2+bx+c$ 配方
& $a\neq0$ & 符号写错；忘记 $a\neq0$ & 一级 \\
\addlinespace
判别式
$\Delta=b^2-4ac$
& 由求根公式中根号内的部分得到
& 方程必须是一元二次 & 只记公式，不会判断根的个数 & 一级 \\
\addlinespace
韦达定理
$x_1+x_2=-\dfrac{b}{a}$，
$x_1x_2=\dfrac{c}{a}$
& 由求根公式相加、相乘得到
& $\Delta\geq0$ & 忽略 $\Delta\geq0$ 的前提 & 二级 \\
\addlinespace
大于取两边，小于取中间
& 由图象在 $x$ 轴上方 / 下方得到
& $a>0$ 且已因式分解 & $a<0$ 时不化正就用，方向错 & 二级 \\
\addlinespace
左同右异
（对称轴在 $y$ 轴左，$a,b$ 同号；在右，异号）
& 由 $x=-\dfrac{b}{2a}$ 与 $0$ 比较得到
& $a\neq0$ & 只记口诀，不会推导 & 三级 \\
\addlinespace
恒成立条件
$ax^2+bx+c>0$ 恒成立
$\iff a>0$ 且 $\Delta<0$
& 由「图象全部在 $x$ 轴上方」得到
& 二次项系数含参时需讨论 & 漏掉 $a=0$ 的讨论 & 二级 \\
\addlinespace
根的分布
（两根都在区间内）
& 用对称轴位置、端点函数值符号、$\Delta$ 联立
& 需分类讨论 & 条件列不全，漏掉 $\Delta$ & 三级 \\
\bottomrule
\end{tabular}
\caption{常见二级结论卡片}
\end{table}

\section{使用建议}

\begin{itemize}
  \item \key{一级}：必须能独立推导，且能讲清楚每一步的依据；
  \item \key{二级}：记住形式，会用；知道它是从哪一层推出来的；
  \item \key{三级}：看懂即可，不要求默写，遇到时知道有这个工具。
\end{itemize}

\section{复习节奏}

同一条结论至少要在三个地方出现：

\begin{enumerate}
  \item 新授课：推导它（第一次）；
  \item 基础练习：直接套用它（第二次）；
  \item 综合题：在变化的条件下辨认它（第三次）。
\end{enumerate}
